\begin{frame}{Marking global inconsistencies?}
  Consider this program:
  \[%
    \ELam{f}{\TUnknown}{\EAp{f}{(\EPlus{f}{1})}}
  \]%

  $\assignType{f}{\TUnknown}$, so the bidirectional system operates gradually, \pause
  but $f$ is 
  %
  \begin{enumerate}
    \item applied as a function \textcolor{Black!50}{(a function?)}
      \pause
    \item an operand of $\ECOpPlus$ \textcolor{Black!50}{(a number?)}
  \end{enumerate}

  \pause
  \vspace{1em}
  \emph{Global, constraint-based type checking would have caught this!}

  \note[item]{Since the marked lambda calculus is for bidirectional systems, it might not do
    everything we'd like}
  \note[item]{In this program, the type of f is given to be unknown, so the bidirectional system
    would operate gradually and results in no marks}
  \note[item]{But we notice that it's used a function and a number simultaneously, which isn't
    possible}
  \note[item]{In other ways, there's no way to replace the unknown type, or type hole as well we
    call them, with a known type}
  \note[item]{And anyway we think, well, some sort of global, constraint-based inference system
    would have caught this!}
\end{frame}

\begin{frame}{Layers upon layers}
  Get the best of both worlds by layering \textbf{constraint-based inference} 
  atop the marked lambda calculus \\[1em]

  \pause
  \begin{itemize}
    \item The marked lambda calculus localizes and recovers predictably from \emph{local inconsistencies}

      \pause
    \item \textbf{Type hole inference} solves and marks \emph{global inconsistencies} 

      \pause
      \begin{itemize}
        \item Downstream service to supplement the marked lambda calculus
      \end{itemize}
  \end{itemize}

  \note[item]{In the second part of this paper, we describe how we can the best of both worlds by
    layering a constraint-based inference system atop the marked lambda calculus}
  \note[item]{Where the latter provides local inference with error localization and recovery}
  \note[item]{And type hole inference operates on the marked terms to solve and mark
    global inconsistencies like the one in the example}
  \note[item]{Rather than being a standalone type checker, this is an example of a downstream
    service that uses the marked terms that result from marking and supplements the marked lambda
    calculus}
  \note[item]{Without getting into the weeds of how exactly it all works---you can refer to the
    paper for the details---let's just take the previous example into our implementation of type
    hole inference in Hazel}
\end{frame}

\newcommand{\Id}[1]{\ensuremath{\goodcolor{hole}{#1}}}
\newcommand{\IdU}{\ensuremath{\Id{u}}}
\newcommand{\ProvU}{\ensuremath{\IdU}}
\newcommand{\ProvExp}[1]{\ensuremath{exp\goodcolor{\colorOkSideJudge}{(}#1\goodcolor{\colorOkSideJudge}{)}}}
\newcommand{\ProvIn}[1]{\ensuremath{\to_L\goodcolor{\colorOkSideJudge}{(}#1\goodcolor{\colorOkSideJudge}{)}}}
\newcommand{\ProvOut}[1]{\ensuremath{\to_R\goodcolor{\colorOkSideJudge}{(}#1\goodcolor{\colorOkSideJudge}{)}}}
\newcommand{\PTS}{\operatorname{\mathsf{PotentialTypeSet}}}
\newcommand{\cursor}{\ensuremath{\goodcolor{\colorOkSideJudge}{\bm{\wedge}}}}

\begin{frame}[fragile, t]{Type hole inference in Hazel}
  \begin{center}
    \alt<6->{
      \alt<7->{
        \includegraphics[width=0.9\textwidth]{inference-arrow.png}
      }{
        \includegraphics[width=0.9\textwidth]{inference-int.png}
      }
    }{
      \includegraphics[width=0.9\textwidth]{inference.png}
    }
    \vspace{2em}
    \alt<6->{
      \alt<7->{
        \includegraphics[width=0.9\textwidth]{inspector-arrow.png}
      }{
        \includegraphics[width=0.9\textwidth]{inspector-int.png}
      }
    }{
      \includegraphics[width=0.9\textwidth]{inspector.png}
    }
  \end{center}

  \vspace{-2em}
  \begin{enumerate}
    \small
    \setlength\itemsep{1pt}
    \item<2-> Gather constraints, treat occurrences of $\TUnknown$ as unification variables
    \item<3-> When solvable, proceed as normal to find substitution
    \item<4-> When unsolvable, maintain partial solutions: possible type fillings, \\
          which are offered to the user for selection
    \item<5-> Control returned to the bidirectional system upon selection
  \end{enumerate}

  % \begin{center}
    % Try it out at \\
    % \textcolor{MidnightBlue}{[\url{https://hazel.org/build/thi}]}
  % \end{center}

  \begin{tikzpicture}[remember picture,overlay]
    \node[xshift=-4.6cm, yshift=-4.7cm, visible on=<6>] at (current page.north east) { \mousePointer };
    \node[xshift=-2.3cm, yshift=-4.7cm, visible on=<7>] at (current page.north east) { \mousePointer };
  \end{tikzpicture}

  \only<1-5>{
  \note[item]{This is a constraint-based inference system, so we gather constraints and try to solve
    them with a standard unification algorithm}
  \note[item]{When unification is able to solve the constraints, things proceed as normal to find
    the appropriate substitution}
  \note[item]{However, like here, when unsolvable, the system degrades gracefully to maintain
    partial solutions by accumulating the possible type fillings, which are then offered to the user
    for selection}
  \note[item]{So notice that the system has marked the unknown type, the type hole here as unsolvable}
  \note[item]{and the cursor inspector tells us that the constraints demand f be both an int and int
    -> unknown}
  \note[item]{Instead of failing outright, type hole inference collects and offers both possible fillings to
    the user, and if we hover over each, we can preview and then click to select the one we intend}
  }
  \only<6->{
  \note[item]{When hovering over Int, it fills Int into the type hole, and now the application of f
    is marked with an error because control has been returned naturally to the bidirectional system}
  \note[item]{Similarly, hovering over Int -> ? results in the second f being marked with an error}
  }
\end{frame}

\begin{frame}
  This approach is \emph{neutral} \\[1em]

  \pause
  \begin{itemize}
    \item Localize errors to the originating type hole or error mark, \\
      instead of guessing about user intent

      \pause
    \item All potential type hole fillings are provided to the user for selection

      \pause
    \item Control returned to bidirectional system after user selects
  \end{itemize}

  \note[item]{This approach is neutral in the sense that we don't try to guess about user intent but
    simply provide all potential type hole fillings to the user for selection}
  \note[item]{facilitating a workflow of filling in type holes and returning control to the marked
    lambda calculus}
  \note[item]{It turns out that the bidirectional and unification approaches can work nicely
    together!}
\end{frame}
